
\documentclass[11pt,a4paper]{article}
\usepackage[T1]{fontenc}
\usepackage{lmodern}
\usepackage{amsmath,amssymb,amsfonts,mathtools}
\usepackage{graphicx,booktabs,longtable,array,tabularx,adjustbox}
\usepackage{geometry,microtype,fancyhdr,lastpage,etoolbox,enumitem}
\usepackage{newunicodechar,textcomp}
\usepackage{hyperref}
\hypersetup{hidelinks,pdfcreator={XeLaTeX / QM1 academic edition}}
\geometry{left=20mm,right=20mm,top=22mm,bottom=23mm}
\setcounter{secnumdepth}{2}
\setlength{\parindent}{0pt}
\setlength{\parskip}{4pt plus 1pt}
\setlength{\emergencystretch}{2em}
\setlength{\tabcolsep}{3pt}
\renewcommand{\arraystretch}{1.16}
\setlist{nosep,topsep=3pt,itemsep=2pt,leftmargin=7mm}
\AtBeginEnvironment{longtable}{\footnotesize\setlength{\tabcolsep}{2.2pt}}
\newcommand{\tightlist}{\setlength{\itemsep}{0pt}\setlength{\parskip}{0pt}}
\providecommand{\Cov}{\operatorname{Cov}}
\providecommand{\Var}{\operatorname{Var}}
\providecommand{\Corr}{\operatorname{Corr}}
\providecommand{\MAD}{\operatorname{MAD}}
\providecommand{\E}{\mathbb{E}}
\providecommand{\R}{\mathbb{R}}
\newunicodechar{—}{\textemdash}
\newunicodechar{–}{\textendash}
\newunicodechar{€}{\texteuro}
\newunicodechar{≤}{\ensuremath{\leq}}
\newunicodechar{≥}{\ensuremath{\geq}}
\newunicodechar{×}{\ensuremath{\times}}
\newunicodechar{÷}{\ensuremath{\div}}
\newunicodechar{Σ}{\ensuremath{\Sigma}}
\newunicodechar{σ}{\ensuremath{\sigma}}
\newunicodechar{μ}{\ensuremath{\mu}}
\newunicodechar{Δ}{\ensuremath{\Delta}}
\newunicodechar{∑}{\ensuremath{\sum}}
\newunicodechar{≈}{\ensuremath{\approx}}
\newunicodechar{√}{\ensuremath{\surd}}
\newunicodechar{≠}{\ensuremath{\ne}}
\newunicodechar{→}{\ensuremath{\to}}
\newunicodechar{∞}{\ensuremath{\infty}}
\newunicodechar{∈}{\ensuremath{\in}}
\newunicodechar{−}{\ensuremath{-}}
\newunicodechar{²}{\ensuremath{^{2}}}
\newunicodechar{³}{\ensuremath{^{3}}}
\newunicodechar{₀}{\ensuremath{_{0}}}
\newunicodechar{₁}{\ensuremath{_{1}}}
\newunicodechar{₂}{\ensuremath{_{2}}}
\newunicodechar{₃}{\ensuremath{_{3}}}
\newunicodechar{₄}{\ensuremath{_{4}}}
\newunicodechar{ₜ}{\ensuremath{_{t}}}
\newunicodechar{ₚ}{\ensuremath{_{p}}}
\newunicodechar{ₙ}{\ensuremath{_{n}}}
\newunicodechar{ₓ}{\ensuremath{_{x}}}
\newunicodechar{ᵢ}{\ensuremath{_{i}}}
\newunicodechar{ŷ}{\ensuremath{\hat y}}
\newunicodechar{Ŷ}{\ensuremath{\widehat Y}}
\newunicodechar{ₘ}{\ensuremath{_{m}}}
\newunicodechar{ₑ}{\ensuremath{_{e}}}
\newunicodechar{ₐ}{\ensuremath{_{a}}}
\newunicodechar{ₖ}{\ensuremath{_{k}}}
\newunicodechar{ₛ}{\ensuremath{_{s}}}
\newunicodechar{ₗ}{\ensuremath{_{l}}}
\newunicodechar{ₕ}{\ensuremath{_{h}}}
\newunicodechar{ₒ}{\ensuremath{_{o}}}
\newunicodechar{₊}{\ensuremath{_{+}}}
\newunicodechar{₋}{\ensuremath{_{-}}}
\newunicodechar{₌}{\ensuremath{_{=}}}
\newunicodechar{₍}{\ensuremath{_{(}}}
\newunicodechar{₎}{\ensuremath{_{)}}}
\newunicodechar{⁻}{\ensuremath{^{-}}}
\newunicodechar{⁺}{\ensuremath{^{+}}}
\newunicodechar{⁰}{\ensuremath{^{0}}}
\newunicodechar{¹}{\ensuremath{^{1}}}
\newunicodechar{⁴}{\ensuremath{^{4}}}
\newunicodechar{⁵}{\ensuremath{^{5}}}
\newunicodechar{⁶}{\ensuremath{^{6}}}
\newunicodechar{⁷}{\ensuremath{^{7}}}
\newunicodechar{⁸}{\ensuremath{^{8}}}
\newunicodechar{⁹}{\ensuremath{^{9}}}
\makeatletter
\renewcommand{\section}{\@startsection{section}{1}{0pt}{2.0ex plus .7ex minus .3ex}{1.0ex}{\large\bfseries}}
\renewcommand{\subsection}{\@startsection{subsection}{2}{0pt}{1.8ex plus .5ex minus .2ex}{.7ex}{\normalsize\bfseries}}
\renewcommand{\subsubsection}{\@startsection{subsubsection}{3}{0pt}{1.3ex plus .4ex minus .2ex}{.5ex}{\normalsize\bfseries}}
\makeatother
\pagestyle{fancy}
\fancyhf{}
\renewcommand{\headrulewidth}{0.35pt}
\renewcommand{\footrulewidth}{0pt}
\fancyhead[L]{\small\textsc{KEDGE Business School --- QM1}}
\fancyhead[R]{\small\textsc{2026--2027}}
\fancyfoot[L]{\small Badr ABOUFARASSE, PhD}
\fancyfoot[R]{\small\thepage\ /\ \pageref{LastPage}}
\setlength{\headheight}{14pt}
\setkeys{Gin}{keepaspectratio}
\begin{document}

\begin{center}
{\small\textsc{KEDGE Business School \quad EBP1 \quad Quantitative Methods 1}}\\[6pt]
{\Large\bfseries Worksheet 6 — Index numbers}\\[6pt]
{\normalsize\bfseries Exercise statements}\\[6pt]
{\small Badr ABOUFARASSE, PhD \quad --- \quad 2026--2027}
\end{center}
\vspace{2mm}\hrule\vspace{6mm}

\hypertarget{core-applications}{%
\section{Core applications}\label{core-applications}}

Master the essential definitions and calculations from the course.

\emph{Core application --- 20 min}

\hypertarget{exercise-1-elementary-indices-rebasing-and-properties}{%
\subsection{Exercise 1 -- Elementary indices, rebasing and properties}\label{exercise-1-elementary-indices-rebasing-and-properties}}

A service costs 80 euros in 2023, 92 euros in 2024 and 101.20 euros in 2025. Calculate the three base-100 indices, verify circularity, calculate the reverse index and verify reversibility, rebase the series to 2024=100, and interpret index 126.5 correctly.

\emph{Core application --- 25 min}

\hypertarget{extra-1-complete-pricequantity-decomposition}{%
\subsection{Extra 1 --- Complete price--quantity decomposition}\label{extra-1-complete-pricequantity-decomposition}}

A purchasing department observes:

\begin{longtable}[]{@{}lllll@{}}
\toprule\noalign{}
Product & p0 & q0 & p1 & q1 \\
\midrule\noalign{}
\endhead
\bottomrule\noalign{}
\endlastfoot
A & 4 & 100 & 5 & 90 \\
B & 10 & 50 & 11 & 60 \\
C & 20 & 20 & 18 & 30 \\
\end{longtable}

\begin{enumerate}
\tightlist
\item
  Calculate the four totals Σp0q0, Σp1q0, Σp0q1 and Σp1q1.
\item
  Calculate the value, Laspeyres-price, Paasche-price and Fisher-price indices.
\item
  Interpret the gap between value and price change.
\end{enumerate}

\textbf{Method hint:}

Complete all four cross-value columns before forming any index.

\hypertarget{assessment-and-exam-practice}{%
\section{Assessment and exam practice}\label{assessment-and-exam-practice}}

Practise independently with assessment and final-exam style questions.

\emph{Assessment / exam practice --- 35 min}

\hypertarget{exercise-2-laspeyres-paasche-fisher-and-factor-reversal}{%
\subsection{Exercise 2 -- Laspeyres, Paasche, Fisher and factor reversal}\label{exercise-2-laspeyres-paasche-fisher-and-factor-reversal}}

\begin{longtable}[]{@{}lllll@{}}
\toprule\noalign{}
& {\(p_0\)} & {\(q_0\)} & {\(p_1\)} & {\(q_1\)} \\
\midrule\noalign{}
\endhead
\bottomrule\noalign{}
\endlastfoot
Product X & 4 & 100 & 5 & 90 \\
Product Y & 10 & 40 & 12 & 50 \\
\end{longtable}

Calculate the four aggregates, Laspeyres, Paasche and Fisher price indices, the corresponding quantity indices, the value index, and verify Fisher factor reversal on a base-100 scale.

\emph{Assessment / exam practice --- 45 min}

\hypertarget{exercise-4-complete-four-product-indices-exam-style}{%
\subsection{Exercise 4 --- Complete four-product indices --- exam style}\label{exercise-4-complete-four-product-indices-exam-style}}

For four products, prices and quantities are observed in two periods. Base 0 = 100.

\begin{longtable}[]{@{}lllll@{}}
\toprule\noalign{}
Produit / Product & p₀ & q₀ & p₁ & q₁ \\
\midrule\noalign{}
\endhead
\bottomrule\noalign{}
\endlastfoot
A & 53 & 20 & 62 & 25 \\
B & 41 & 28 & 52 & 34 \\
C & 75 & 35 & 82 & 38 \\
D & 25 & 30 & 38 & 33 \\
\end{longtable}

\begin{enumerate}
\tightlist
\item
  Build the four product columns p₀q₀, p₁q₀, p₀q₁, p₁q₁ and totals.
\item
  Compute Laspeyres, Paasche and Fisher price and quantity indices.
\item
  Compute the value index and verify Fisher\textquotesingle s factor reversal. Interpret the three indices.
\end{enumerate}

\hypertarget{optional-advanced-practice}{%
\section{Optional advanced practice}\label{optional-advanced-practice}}

Explore additional interpretations and professional applications.

\emph{Advanced practice --- 15 min}

\hypertarget{exercise-3-chain-indices-and-deflation}{%
\subsection{Exercise 3 -- Chain indices and deflation}\label{exercise-3-chain-indices-and-deflation}}

Annual link indices are 105, 108 and 97. Calculate the cumulative chain index and interpret it. Explain why annual changes must be compounded rather than added. Then deflate nominal revenue of 2.4 million euros when the price index equals 120.

\emph{Advanced practice --- 25 min}

\hypertarget{extra-2-rebasing-inflation-and-real-performance}{%
\subsection{Extra 2 --- Rebasing, inflation and real performance}\label{extra-2-rebasing-inflation-and-real-performance}}

A nominal revenue series is €500k, €560k and €630k. The corresponding price index (year 1 = 100) is 100, 108 and 120.

\begin{enumerate}
\tightlist
\item
  Calculate real revenue in year-1 euros.
\item
  Calculate nominal and real growth from year 1 to year 3.
\item
  Rebase the price index so year 2 = 100.
\item
  Explain the result to a manager who claims that activity grew by 26\%.
\end{enumerate}

\textbf{Method hint:}

Deflate by dividing nominal value by index/100. Rebasing changes the reference, not the underlying inflation path.

\end{document}
